% Figure from MATH 590 notes, section 16. Compile with the notes preamble.
\begin{tikzpicture}[scale=1.35]
  % one cover element highlighted
  \fill[figB!10] (1.3,0) circle (0.8);
  % the space X
  \draw[figB, thick] plot[smooth cycle, tension=0.8] coordinates
    {(-1.7,0.2) (-1.3,0.85) (-0.5,0.8) (0.2,1.2) (0.8,0.8) (1.8,0.4) (1.9,-0.3) (1.0,-0.5) (0.4,-1.2) (-0.4,-1.0) (-1.2,-0.6)};
  \node[font=\small, figB] at (-1.85,-0.55) {$X$};
  % the open cover
  \draw[black!55, dashed] (-1,0.2) circle (1.0) (0.3,0.6) circle (0.9) (0.2,-0.5) circle (0.9);
  \draw[figB, dashed, thick] (1.3,0) circle (0.8);
  \node[font=\scriptsize, black!65] at (-1.75,1.05) {$U_1$};
  \node[font=\scriptsize, black!65] at (0.05,1.62) {$U_2$};
  \node[font=\scriptsize, black!65] at (-0.55,-1.35) {$U_3$};
  \node[font=\scriptsize, figB] at (2.2,-0.6) {$U_4$};
  % small sets of diameter < delta: each sits inside one U
  \foreach \p in {(-1.35,0.1), (-0.45,0.35), (0.75,0.0), (1.5,0.1), (0.1,-0.85), (0.35,0.85)}
    \filldraw[figR, fill=figR!25, thick] \p circle (0.17);
  % scale bar
  \draw[figR, thick, |-|] (-0.3,-1.72) -- (0.28,-1.72);
  \node[font=\scriptsize, figR, right] at (0.32,-1.72) {$\delta$};
\end{tikzpicture}
